\documentclass[a4paper,12pt]{article}
	\usepackage[ansinew]{inputenc}
	\usepackage{ae}	
	\usepackage[brazil]{babel}
	\usepackage{xcolor}
	\usepackage{graphicx}
	
	% O pacote pinlabel permite inserir textos em figuras EPS e PDF
	\usepackage{pinlabel}
	
	\newcommand{\comando}[1]{{\normalfont\textbackslash\color{blue!50!black}#1}}
\newcommand{\pacote}[1]{{\normalfont\sffamily\color{orange!50!black}#1}}
\newcommand{\ambiente}[1]{{\normalfont\sffamily\color{green!50!black}#1}}
\newcommand{\contador}[1]{{\normalfont\sffamily\color{red!50!black}#1}}
\newcommand{\foreign}[1]{\textsl{#1}}
\newcommand\linhabase{\noindent\makebox[0pt][l]{\color{red}\rule{\textwidth}{0.1pt}}}
	
	\graphicspath{
		{./figuras/EPS/}
		{./figuras/JPG/}
		{./figuras/PDF/}
		{./figuras/PNG/}
		}
	
\begin{document}

	\begin{figure}		
		\labellist
		\small
		\hair 1pt
		\pinlabel {Obrigado, \pacote{pinlabel}!} [lb] at 103 255
		%\pinlabel {Obrigado, \pacote{pinlabel}!} [B] at 215 425
		\endlabellist
		\centering
		\includegraphics[width=10cm]{Tei-Gi}
		\caption{diferentemente do \pacote{psfrag}, que \emph{substitui} textos existentes no arquivo EPS por outros, o pacote \pacote{pinlabel} \emph{insere-os} numa posição qualquer da figura, seja ela EPS ou PDF (a versão EPS é sempre necessária). As coordenadas dos rótulos, em \texttt{pt}, podem ser facilmente determinadas abrindo-se o arquivo EPS no Ghostview. Uma observação importante: o nome do arquivo inserido através de \comando{includegraphics} deve ser dado \emph{sem extensão}. Além disso, ao menos o arquivo EPS deve estar na mesma pasta do arquivo \LaTeX, haja vista que o pacote \pacote{pinlabel} não reconhece a instrução \comando{graphicspath}.}
	\end{figure}
	
	% EXERCÍCIO: abra o arquivo Tei-Gi.eps e reposicione o texto do exemplo acima no centro do "olho" esquerdo, ou acrescente outro texto nesta posição.

\end{document}
